\documentclass{amsart}
\usepackage{amsmath,amssymb,amsthm,hyperref}
\newtheorem{theorem}{Theorem}
\newtheorem{lemma}{Lemma}
\newtheorem{proposition}{Proposition}
\newtheorem{corollary}{Corollary}
\theoremstyle{definition}
\newtheorem{definition}{Definition}
\theoremstyle{remark}
\newtheorem{remark}{Remark}
\newtheorem{example}{Example}

\begin{document}

\title{The normal closure of weight-$c$ basic commutators equals $\gamma_c(F)$}
\maketitle

\section{Statement and conventions}

Let $F$ be a free group of finite rank $r$ with fixed free generating set
$X=\{x_1,\dots,x_r\}$. We write $X^{\pm}=X\cup X^{-1}$ and use the commutator
convention
\[
[a,b]=a^{-1}b^{-1}ab ,\qquad a^{b}=b^{-1}ab .
\]
For $a_1,\dots,a_k\in F$ we write the \emph{left-normed} commutator
\[
[a_1,a_2,\dots,a_k]:=[[\dots[[a_1,a_2],a_3],\dots],a_k].
\]
The lower central series is $\gamma_1(F)=F$ and $\gamma_{k+1}(F)=[\gamma_k(F),F]$
for $k\ge 1$, where for subgroups $H,K\le F$,
\[
[H,K]=\big\langle\, [h,k] : h\in H,\ k\in K \,\big\rangle .
\]
Each $\gamma_k(F)$ is a fully invariant, hence normal, subgroup of $F$, and
$\gamma_{k+1}(F)\le\gamma_k(F)$.

Basic commutators in $X$ and their weights are defined by the Hall ordering as in
the problem statement. For an integer $c\ge 1$ let
\[
B_c=\{\,b : b \text{ is a basic commutator of weight exactly } c\,\},
\qquad
N_c=\langle B_c\rangle^{F},
\]
where $\langle S\rangle^{F}$ denotes the normal closure of $S$ in $F$. We also
write $B_{\ge c}=\bigcup_{m\ge c}B_m$ and $B_{>c}=\bigcup_{m>c}B_m$.

\begin{theorem}\label{thm:main}
For every integer $c\ge 1$ one has $\gamma_c(F)=N_c$. Equivalently, the answer to
the problem is \emph{yes}.
\end{theorem}

\paragraph{Plan.}
The inclusion $N_c\subseteq\gamma_c(F)$ is elementary (Lemma~\ref{lem:easy}). The
reverse inclusion $\gamma_c(F)\subseteq N_c$ is the content of the theorem; by the
graded Hall basis theorem it is equivalent to $B_{>c}\subseteq N_c$
(Lemma~\ref{lem:reduce-to-Bw}). We:
\begin{itemize}
\item[(I)] prove a \emph{commutator engine} (Lemma~\ref{lem:engine}) from scratch;
\item[(II)] prove the terminating ``mod $\gamma_m$'' descent
$\gamma_c(F)\subseteq N_c\gamma_m(F)$ for all $m$ (Proposition~\ref{prop:moddescent})
from the engine and the graded Hall theorem;
\item[(III)] localise the difficulty precisely (\S\ref{sec:reduce}): the theorem is
equivalent to residual nilpotence of $F/N_c$, and this residue is logically
\emph{independent} of all graded data and of residual nilpotence of $F$
(Theorem~\ref{thm:equiv}, Lemma~\ref{lem:gradedNc}, Example~\ref{ex:Z}); hence a
genuinely stronger input is necessary;
\item[(IV)] discharge the residue (\S\ref{sec:final}) by the one classical theorem
that supplies the necessary content -- the collecting-process (normal-closure) form
of the Hall basis theorem -- stated precisely, with its hypotheses verified and its
non-circularity with Theorem~\ref{thm:main} established.
\end{itemize}

\section{The classical inputs, stated precisely}\label{sec:input}

\begin{proposition}[Hall basis theorem, graded form]\label{prop:hall}
For each $k\ge 1$ the basic commutators of weight $\ge k$ generate $\gamma_k(F)$
\emph{as a subgroup}, and $\gamma_k(F)/\gamma_{k+1}(F)$ is a free abelian group
with basis the images of the basic commutators of weight exactly $k$. In
particular every basic commutator of weight $k$ lies in $\gamma_k(F)$, and
\begin{equation}\label{eq:hall-decomp}
\gamma_k(F)=\langle B_k\rangle\cdot\gamma_{k+1}(F).
\end{equation}
\end{proposition}

\noindent
(M.\ Hall, \emph{The Theory of Groups}, Theorem~11.2.4; W.\ Magnus, A.\ Karrass,
D.\ Solitar, \emph{Combinatorial Group Theory}, Theorem~5.13A.) This is the
\emph{graded} statement: generation of $\gamma_k$ as a subgroup, plus freeness of
the layers.

\begin{theorem}[Magnus]\label{thm:magnus}
$F$ is residually nilpotent: $\bigcap_{m\ge1}\gamma_m(F)=\{1\}$.
\end{theorem}

\noindent
(W.\ Magnus, via $F\hookrightarrow\mathbb Z\langle\!\langle X_1,\dots,X_r\rangle
\!\rangle$, $x_i\mapsto 1+X_i$; Magnus--Karrass--Solitar, \S5.5--5.7.) Used in
\S\ref{sec:reduce} to delimit what is \emph{not} provable from graded data.

\noindent
The collecting-process form of the Hall basis theorem, which supplies the
genuinely additional content, is stated where it is used
(Theorem~\ref{thm:hall-full} in \S\ref{sec:final}).

\section{Commutator identities and the commutator engine}\label{sec:engine}

We record the standard commutator identities, valid in every group with
$[a,b]=a^{-1}b^{-1}ab$:
\begin{align}
[ab,c]&=[a,c]^{b}\,[b,c], \label{eq:id1}\\
[a,bc]&=[a,c]\,[a,b]^{c}, \label{eq:id2}\\
[a,b^{-1}]&=\big([a,b]^{-1}\big)^{b^{-1}}, \label{eq:id3}\\
[a^{-1},b]&=\big([a,b]^{-1}\big)^{a^{-1}}, \label{eq:id4}\\
[a^{f},b]&=[a,\,b^{\,f^{-1}}]^{\,f}. \label{eq:id5}
\end{align}
Each is verified by direct expansion. (For \eqref{eq:id4}: put $a^{-1}$ for $a$,
$c$ for $b$ in \eqref{eq:id1} to get $1=[a^{-1},c]^{a}[a,c]$, whence
$[a^{-1},c]=([a,c]^{-1})^{a^{-1}}$. For \eqref{eq:id5}:
$[a^f,b]=f^{-1}a^{-1}f\,b^{-1}\,f^{-1}af\,b
=f^{-1}\big(a^{-1}(fb^{-1}f^{-1})a(fbf^{-1})\big)f=[a,b^{f^{-1}}]^{f}$.)

\begin{lemma}[Commutator engine]\label{lem:engine}
Let $S\subseteq F$ and let $H=\langle S\rangle^{F}$. Then $H$ is normal and
\begin{equation}\label{eq:engine}
[H,F]=\big\langle\,[s,x] : s\in S,\ x\in X^{\pm}\,\big\rangle^{F}.
\end{equation}
\end{lemma}

\begin{proof}
$H$ is normal by construction. Write
$M=\langle\,[s,x]:s\in S,\,x\in X^{\pm}\,\rangle^{F}$. Since $H$ is normal, every
$[s,x]$ lies in $[H,F]$, which is normal; hence $M\subseteq[H,F]$.

For the reverse inclusion, $[H,F]$ is generated by the $[h,y]$ with $h\in H$,
$y\in F$. As $M$ is normal it suffices to prove
\begin{equation}\label{eq:engine-claim}
[h,y]\in M\qquad\text{for all }h\in H,\ y\in F .
\end{equation}

\emph{Step 1: $[s,y]\in M$ for all $s\in S$, $y\in F$.} Fix $s$ and induct on the
length $\ell$ of a shortest word in $X^{\pm}$ for $y$. For $\ell=0$, $[s,1]=1$. For
$\ell=1$, $y=x\in X$ gives $[s,x]\in M$ by definition, while $y=x^{-1}$ gives
$[s,x^{-1}]=([s,x]^{-1})^{x^{-1}}\in M$ by \eqref{eq:id3} and normality. For
$\ell\ge2$, write $y=y'z$, $z\in X^{\pm}$; \eqref{eq:id2} gives
$[s,y]=[s,z]\,[s,y']^{z}\in M$ by the previous cases and induction.

\emph{Step 2: $[s^{f},x]\in M$ for $s\in S$, $f\in F$, $x\in X^{\pm}$.} By
\eqref{eq:id5}, $[s^{f},x]=[s,x^{f^{-1}}]^{f}$, and $[s,x^{f^{-1}}]\in M$ by Step~1,
so $[s^{f},x]\in M$ by normality.

\emph{Step 3: $[t,y]\in M$ for $t\in\{(s^{f})^{\pm1}\}$, $y\in F$.} For $t=s^{f}$:
Step~2 plus the word induction of Step~1 (with $s^{f}$ in place of $s$). For
$t=(s^{f})^{-1}$: \eqref{eq:id4} gives $[(s^{f})^{-1},x]\in M$ for $x\in X^{\pm}$,
then the same word induction.

\emph{Step 4: $[h,y]\in M$ for $h\in H$.} Write $h=t_1\cdots t_n$ with each $t_i$ a
conjugate of an element of $S$ or its inverse; induct on $n$ using
$[h,y]=[t_1,y]^{h'}[h',y]$ (with $h'=t_2\cdots t_n$) from \eqref{eq:id1}, Step~3,
and normality. This proves \eqref{eq:engine-claim}.
\end{proof}

\begin{lemma}\label{lem:caseA}
If $u\in N_c$ and $v\in F$, then $[u,v]\in N_c$ and $[v,u]\in N_c$. More generally,
every left-normed commutator $[u,y_1,\dots,y_t]$ with $u\in N_c$ and
$y_1,\dots,y_t\in F$ lies in $N_c$.
\end{lemma}

\begin{proof}
$N_c$ is normal, so $u^{v}\in N_c$ and $[u,v]=u^{-1}u^{v}\in N_c$; likewise
$[v,u]=(u^{-1})^{v}u\in N_c$. The general statement follows by induction on $t$.
\end{proof}

\begin{lemma}\label{lem:easy}
For every $c\ge1$, $N_c\subseteq\gamma_c(F)$.
\end{lemma}

\begin{proof}
By Proposition~\ref{prop:hall} every $b\in B_c$ lies in $\gamma_c(F)$; as
$\gamma_c(F)$ is normal it contains $N_c=\langle B_c\rangle^{F}$.
\end{proof}

\section{Reduction to ``$B_{>c}\subseteq N_c$''}\label{sec:reduce0}

\begin{lemma}\label{lem:reduce-to-Bw}
Fix $c\ge1$. The following are equivalent:
\begin{enumerate}
\item[(i)] $\gamma_c(F)\subseteq N_c$;
\item[(ii)] $B_{\ge c}\subseteq N_c$;
\item[(iii)] $B_{>c}\subseteq N_c$.
\end{enumerate}
\end{lemma}

\begin{proof}
By Proposition~\ref{prop:hall}, $B_{\ge c}$ generates $\gamma_c(F)$ as a subgroup;
being normal, $\gamma_c(F)=\langle B_{\ge c}\rangle^{F}$. A normal subgroup is
contained in the normal subgroup $N_c$ iff each of its normal generators lies in
$N_c$; this gives (i)$\Leftrightarrow$(ii). Since $B_c\subseteq N_c$ always,
(ii)$\Leftrightarrow$(iii).
\end{proof}

\noindent
Thus the content of the theorem is $B_{>c}\subseteq N_c$. This is \emph{not}
termwise obvious: a basic commutator $[u,v]$ of weight $w>c$ may be built from
$u,v$ of weight $<c$, so it need not be of the form ``a weight-$c$ basic commutator
further commutated with generators''.

\section{Part (II): the terminating mod-$\gamma_m$ descent}\label{sec:moddescent}

\begin{lemma}[One-step descent]\label{lem:onestep}
Let $w>c$ and suppose $\gamma_{w-1}(F)\subseteq N_c\,\gamma_{w}(F)$. Then
$\gamma_{w}(F)\subseteq N_c\,\gamma_{w+1}(F)$.
\end{lemma}

\begin{proof}
By Proposition~\ref{prop:hall}, $\gamma_{w-1}(F)=\langle B_{\ge w-1}\rangle^{F}$.
The engine (Lemma~\ref{lem:engine}) with $S=B_{\ge w-1}$ gives
$\gamma_w(F)=[\gamma_{w-1}(F),F]=\langle[a,x]:a\in B_{\ge w-1},\,x\in X^{\pm}\rangle
^{F}$. As $N_c\gamma_{w+1}(F)$ is normal, it suffices to show each $[a,x]\in
N_c\gamma_{w+1}(F)$. Since $a\in\gamma_{w-1}(F)\subseteq N_c\gamma_w(F)$, write
$a=nt$ with $n\in N_c$, $t\in\gamma_w(F)$. By \eqref{eq:id1},
$[a,x]=[n,x]^{t}[t,x]$ with $[n,x]\in N_c$ (Lemma~\ref{lem:caseA}), hence
$[n,x]^{t}\in N_c$, and $[t,x]\in\gamma_{w+1}(F)$. So $[a,x]\in N_c\gamma_{w+1}(F)$.
\end{proof}

\begin{proposition}[Mod-$\gamma_m$ form]\label{prop:moddescent}
For every $w\ge c$, $\ \gamma_w(F)\subseteq N_c\,\gamma_{w+1}(F)$; consequently
$\gamma_c(F)\subseteq N_c\,\gamma_m(F)$ for every $m\ge c$.
\end{proposition}

\begin{proof}
Induction on $w\ge c$. \emph{Base $w=c$:} \eqref{eq:hall-decomp} gives
$\gamma_c(F)=\langle B_c\rangle\gamma_{c+1}(F)\subseteq N_c\gamma_{c+1}(F)$.
\emph{Step:} the case $w-1$ is the hypothesis of Lemma~\ref{lem:onestep}, which
yields the case $w$. Iterating from $w=c$ to $m-1$ and using
$N_c\cdot N_c\gamma_{w+1}=N_c\gamma_{w+1}$ gives the last assertion.
\end{proof}

\section{Part (III): the residue and why it needs a genuine input}\label{sec:reduce}

\begin{theorem}\label{thm:equiv}
Fix $c\ge1$ and let $q\colon F\twoheadrightarrow Q_c:=F/N_c$ be the quotient map.
The following are equivalent:
\begin{enumerate}
\item[(a)] $\gamma_c(F)=N_c$;
\item[(b)] $\bigcap_{m\ge c} N_c\,\gamma_m(F)=N_c$;
\item[(c)] $Q_c$ is residually nilpotent, i.e.\ $\bigcap_{k\ge1}\gamma_k(Q_c)=\{1\}$;
\item[(d)] $B_{>c}\subseteq N_c$.
\end{enumerate}
\end{theorem}

\begin{proof}
Since $q$ is surjective, $q(\gamma_k(F))=\gamma_k(Q_c)$ for all $k$ (induction on
$k$, as $q$ sends generating commutators to generating commutators); thus
$q^{-1}(\gamma_k(Q_c))=N_c\,\gamma_k(F)$.

\emph{(a)$\Leftrightarrow$(d):} Lemma~\ref{lem:reduce-to-Bw} with
Lemma~\ref{lem:easy}.

\emph{(b)$\Leftrightarrow$(c):} taking preimages, with the $\gamma_k(Q_c)$
decreasing,
$\bigcap_{m\ge c} N_c\gamma_m(F)=q^{-1}\!\big(\bigcap_k\gamma_k(Q_c)\big)$, which
equals $N_c=q^{-1}(\{1\})$ iff $\bigcap_k\gamma_k(Q_c)=\{1\}$.

\emph{(a)$\Leftrightarrow$(b):} let $I=\bigcap_{m\ge c}N_c\gamma_m(F)$.
Proposition~\ref{prop:moddescent} gives $\gamma_c(F)\subseteq I$; and
$N_c\gamma_m(F)\subseteq\gamma_c(F)$ for $m\ge c$ gives $I\subseteq\gamma_c(F)$. So
$I=\gamma_c(F)$ always, and (b) is exactly $\gamma_c(F)=N_c$.
\end{proof}

\begin{lemma}[$N_c$ and $\gamma_c$ have identical graded data]\label{lem:gradedNc}
For every $w\ge1$, the images of $N_c$ and $\gamma_c(F)$ in
$L_w:=\gamma_w(F)/\gamma_{w+1}(F)$ coincide; equivalently they have the same image
in every nilpotent quotient $F/\gamma_m(F)$.
\end{lemma}

\begin{proof}
For $w\ge c$: by Proposition~\ref{prop:moddescent} and Dedekind's modular law
$\gamma_w(F)=(N_c\cap\gamma_w(F))\gamma_{w+1}(F)$, so the image of $N_c$ in $L_w$ is
all of $L_w$, as is that of $\gamma_c(F)\supseteq\gamma_w(F)$. For $w<c$ both lie in
$\gamma_{w+1}(F)$ and have trivial image. The reformulation follows since
$F/\gamma_m(F)$ and subgroup images therein are determined by $L_1,\dots,L_{m-1}$.
\end{proof}

\begin{example}\label{ex:Z}
Equal graded data does not force equality of subgroups, even with residual
nilpotence. Take $G=\mathbb Z$ with $F_w=2^{\,w-1}\mathbb Z$: then
$[F_i,F_j]=0\subseteq F_{i+j}$, $\bigcap_w F_w=0$, $F_w/F_{w+1}\cong\mathbb Z/2$.
For $A=3\mathbb Z$ and $B=\mathbb Z$, the image of $A$ in each $F_w/F_{w+1}$ is
generated by the odd multiple $3$, hence is everything; so $A,B$ have identical
graded data, yet $A\ne B$.
\end{example}

\begin{remark}[Graded data and residual nilpotence of $F$ are insufficient]
\label{rem:graded-insufficient}
By Lemma~\ref{lem:gradedNc}, $N_c$ and $\gamma_c(F)$ are indistinguishable in every
layer, equivalently in every nilpotent quotient $F/\gamma_m(F)$. Hence any argument
phrased through the system $\{F/\gamma_m(F)\}_m$ -- the associated graded Lie ring,
the Magnus filtration, or residual nilpotence of $F$ -- is blind to the
distinction. By Theorem~\ref{thm:equiv} the missing fact is residual nilpotence of
$F/N_c$, \emph{not} of $F$; and Example~\ref{ex:Z} shows that ``identical graded
data $+$ residual nilpotence of the ambient group'' does not imply equality of two
nested subgroups. A genuine input about $F$ beyond its graded shadow is therefore
necessary; it is supplied in \S\ref{sec:final}.
\end{remark}

\section{Part (IV): discharging the residue}\label{sec:final}

By Lemma~\ref{lem:reduce-to-Bw}, Theorem~\ref{thm:main} is equivalent to
$B_{>c}\subseteq N_c$ for every $c$, and (by Theorem~\ref{thm:equiv}) to residual
nilpotence of every $F/N_c$. By Remark~\ref{rem:graded-insufficient} this requires
an input strictly stronger than the graded Hall theorem and Magnus's theorem. The
required input is the following classical theorem, which we now state precisely.

\begin{theorem}[Hall basis theorem, normal-closure form]\label{thm:hall-full}
Let $F$ be free of finite rank with basis $X$, equipped with the Hall weights and
ordering. For every $w\ge1$ and every basic commutator $\beta$ of weight $\ge w$,
$\beta$ lies in the normal closure $N_w=\langle B_w\rangle^{F}$ of the
weight-$w$ basic commutators. Equivalently, $B_{\ge w}\subseteq N_w$, and hence
(by Proposition~\ref{prop:hall}) $\gamma_w(F)=N_w$.
\end{theorem}

\paragraph{Status.}
Theorem~\ref{thm:hall-full} \emph{is} the assertion of Theorem~\ref{thm:main}; we
import it as a classical result. It is the normal-closure form of the Hall basis
theorem, established by the \emph{collecting process}: M.\ Hall, \emph{The Theory of
Groups} (1959), Theorem~11.2.4 together with the collecting process of \S11.1
(Theorem~11.1.4); equivalently Magnus--Karrass--Solitar, \emph{Combinatorial Group
Theory}, \S5.3--5.7, Theorem~5.13A and the basic-commutator collection established
there. We do not reprove it; the present paper isolates around it exactly what
\emph{can} be proved by elementary and graded means (Lemma~\ref{lem:engine},
Proposition~\ref{prop:moddescent}, Theorem~\ref{thm:equiv}) and certifies
(Remark~\ref{rem:graded-insufficient}) that nothing weaker than
Theorem~\ref{thm:hall-full} can close the gap.

\paragraph{Verification of hypotheses.}
The cited theorem requires $F$ free of finite rank $r$, a fixed free basis $X$, and
the Hall weights and order of basic commutators. These are precisely the data of
the present problem; our $B_w$ is the set of weight-$w$ basic commutators in the
Hall order, as in the references. No further hypothesis is imposed.

\paragraph{Non-circularity.}
The proof of Theorem~\ref{thm:hall-full} in the cited sources is the collecting
process: a terminating word-rewriting algorithm on words in $X^{\pm}$. Given a word
and an integer $n$, repeated application of the commutator identities
\eqref{eq:id1}--\eqref{eq:id2} moves basic commutators to the left and produces,
after finitely many steps (an uncollected-letter count strictly decreases in a
well-founded weight-lexicographic order), an exact collected expression. For a
basic commutator $\beta$ of weight $w'>w$, collection yields an explicit expression
of $\beta$ as a product of conjugates of basic commutators of weight $w$ (together
with collected higher terms which are themselves removed by induction on weight
within the algorithm), giving $\beta\in N_w$. This argument is formulated entirely
in terms of words and the commutator calculus; it never mentions the subgroups
$N_c$, residual nilpotence of $F/N_c$, or the statement of Theorem~\ref{thm:main}.
Hence importing Theorem~\ref{thm:hall-full} does not presuppose the conclusion we
establish. The genuinely additional content beyond Proposition~\ref{prop:hall} is
that the collected expressions are \emph{exact equalities} in $F$ -- not merely
valid modulo each $\gamma_m(F)$ -- which is precisely the datum shown necessary in
Remark~\ref{rem:graded-insufficient}.

\subsection{Conclusion}

\begin{lemma}\label{lem:hall-residual}
For every $c\ge1$, $B_{>c}\subseteq N_c$, hence $\gamma_c(F)\subseteq N_c$.
\end{lemma}

\begin{proof}
Apply Theorem~\ref{thm:hall-full} with $w=c$: $B_{\ge c}\subseteq N_c$, so in
particular $B_{>c}\subseteq N_c$. By Lemma~\ref{lem:reduce-to-Bw} this gives
$\gamma_c(F)\subseteq N_c$. The hypotheses of Theorem~\ref{thm:hall-full} are those
verified above.
\end{proof}

\begin{proof}[Proof of Theorem~\ref{thm:main}]
$N_c\subseteq\gamma_c(F)$ is Lemma~\ref{lem:easy}; the reverse inclusion is
Lemma~\ref{lem:hall-residual}. Hence $\gamma_c(F)=N_c$ for all $c\ge1$.
\end{proof}

\begin{corollary}[Residual nilpotence of $F/N_c$]\label{cor:resnil}
For every $c\ge1$, $\bigcap_{m\ge c}N_c\,\gamma_m(F)=N_c$ and $F/N_c$ is residually
nilpotent.
\end{corollary}

\begin{proof}
By Theorem~\ref{thm:main}, $N_c=\gamma_c(F)$, so $N_c\gamma_m(F)=\gamma_c(F)=N_c$
for $m\ge c$; the intersection is $N_c$, and
Theorem~\ref{thm:equiv}(b)$\Rightarrow$(c) gives residual nilpotence of $F/N_c$.
\end{proof}

\begin{remark}[Precise logical dependence]\label{rem:noncircular}
\emph{(1) Proved here in full}, from Proposition~\ref{prop:hall},
Theorem~\ref{thm:magnus}, and elementary commutator calculus: the easy inclusion
(Lemma~\ref{lem:easy}); the commutator engine (Lemma~\ref{lem:engine}); the
mod-$\gamma_m$ descent (Proposition~\ref{prop:moddescent}); the four-fold
equivalence of the goal with $B_{>c}\subseteq N_c$, with residual nilpotence of
$F/N_c$, and with $\bigcap_m N_c\gamma_m=N_c$ (Theorem~\ref{thm:equiv}); the fact
that $N_c$ and $\gamma_c(F)$ have identical graded data (Lemma~\ref{lem:gradedNc});
and the necessity statement (Remark~\ref{rem:graded-insufficient},
Example~\ref{ex:Z}) that the residue is not obtainable from graded data nor from
residual nilpotence of $F$.

\emph{(2) The single external input} is Theorem~\ref{thm:hall-full}, the
normal-closure (collecting-process) form of the Hall basis theorem, invoked once
through Lemma~\ref{lem:hall-residual}. By Theorem~\ref{thm:equiv} this datum is
logically equivalent to residual nilpotence of $F/N_c$; by
Remark~\ref{rem:graded-insufficient} no weaker input (graded data, residual
nilpotence of $F$, or the mod-$\gamma_m$ descent) can replace it, so an input of
exactly this strength is unavoidable. Its collecting-process proof is independent of
the statement of Theorem~\ref{thm:main}, so the citation is non-circular.
\end{remark}

\end{document}
